%! Mode:: "TeX:UTF-8"
%! TEX program = xelatex
%
% CUMCM paper skeleton.
%
% Compile with:  latexmk -xelatex -interaction=nonstopmode paper.tex
% Prerequisite:  cumcmthesis.cls must sit next to this file (see README.md).
% Everything below compiles as-is; replace the placeholders with real content.
%
\PassOptionsToPackage{quiet}{xeCJK}
\documentclass[withoutpreface,notoc]{cumcmthesis}

\usepackage{amsmath,amssymb}
\usepackage{booktabs}
\usepackage{graphicx}
\usepackage{tabularx}

\newcolumntype{C}{>{\centering\arraybackslash}X}
\newcolumntype{L}{>{\raggedright\arraybackslash}X}

\title{基于机理分析与优化模型的定日镜场设计}

\begin{document}
\maketitle

\begin{abstract}
  摘要需要概括每个子问题所用的模型、关键算法与主要定量结论，最后写关键词。
  先写完全文结果，再回来完成摘要。

  \textbf{针对问题一，} 建立了\ldots 模型，采用\ldots 求解，得到\ldots。

  \textbf{针对问题二，} 应用了\ldots 算法，发现\ldots。

  \keywords{关键词一\quad 关键词二\quad 关键词三}
\end{abstract}

\section{问题重述}
简要重述赛题，厘清需要解决的关键问题点。

\section{问题分析}
逐个子问题给出分析思路与拆解方式，并说明题目的已知条件与缺口。

\section{模型假设与符号说明}
本文做出以下假设：
\begin{enumerate}
  \item 假设一\ldots
  \item 假设二\ldots
\end{enumerate}

主要符号如\cref{tab:symbols}所示。

\begin{table}[H]
  \centering
  \caption{符号说明表}
  \label{tab:symbols}
  \begin{tabularx}{\textwidth}{CLC}
    \toprule
    符号 & 说明 & 单位 \\
    \midrule
    $m$ & 质量 & kg \\
    $t$ & 时间 & s \\
    \bottomrule
  \end{tabularx}
\end{table}

\section{问题一的模型建立与求解}
\subsection{模型建立}
给出目标函数、决策变量与约束条件，并说明模型选择的理由与局限。
\begin{equation}
  \label{eq:objective}
  \min_{x \in X} f(x) \quad \text{s.t.} \quad g_i(x) \le 0 .
\end{equation}

\subsection{模型求解}
说明算法步骤、参数取值与实现方式，固定随机种子以便复现。

\subsection{结果与分析}
给出求解结果。图表紧随首次引用处，例如\cref{fig:result}。
\begin{figure}[htbp]
  \centering
  % \includegraphics[width=0.7\textwidth]{figures/result.pdf}
  \fbox{\parbox{0.7\textwidth}{\centering 在此插入由 \texttt{code/} 生成的结果图}}
  \caption{结果图示例}
  \label{fig:result}
\end{figure}

\section{模型检验}
给出灵敏度分析、误差指标与与基线方法的对比。

\section{模型评价与推广}
列出模型的优点、缺点与可推广的方向。

\bibliographystyle{gbt7714-numerical}
\bibliography{refs}

\newpage
\begin{appendices}
  \section{支撑材料列表}
  \begin{table}[H]
    \centering
    \caption{附录文件列表}
    \label{tab:files}
    \begin{tabularx}{\textwidth}{LL}
      \toprule
      文件名 & 功能描述 \\
      \midrule
      code/q1.py & 问题一求解程序 \\
      data/raw.xlsx & 原始数据 \\
      \bottomrule
    \end{tabularx}
  \end{table}
\end{appendices}

\end{document}
